\documentclass[11pt,a4paper]{scrartcl}
% =====================================================================
%  pds-apuntes.tex — Estilo común de apuntes y hojas de ejercicios.
%
%  Uso (desde teoría/Tx/apuntes/ o teoría/Tx/ejercicios/):
%     \documentclass[11pt,a4paper]{scrartcl}   % o scrreprt
%     % =====================================================================
%  pds-apuntes.tex — Estilo común de apuntes y hojas de ejercicios.
%
%  Uso (desde teoría/Tx/apuntes/ o teoría/Tx/ejercicios/):
%     \documentclass[11pt,a4paper]{scrartcl}   % o scrreprt
%     % =====================================================================
%  pds-apuntes.tex — Estilo común de apuntes y hojas de ejercicios.
%
%  Uso (desde teoría/Tx/apuntes/ o teoría/Tx/ejercicios/):
%     \documentclass[11pt,a4paper]{scrartcl}   % o scrreprt
%     \input{../../comun/pds-apuntes}
%     \tema{1}{Introducción: señales y sistemas}
%     \begin{document} \cabecera{Apuntes} ... \end{document}
%  Compilar con:  latexmk -xelatex  (o «make» desde teoría/)
% =====================================================================

\input{../../comun/pds-comun.tex}

\usepackage[a4paper, top=2.6cm, bottom=2.6cm, left=2.5cm, right=2.5cm,
            headheight=1.1cm]{geometry}
\usepackage{fontspec}
\usepackage{unicode-math}
\setmainfont{LibertinusSerif}[Extension=.otf, UprightFont=*-Regular, BoldFont=*-Semibold,
  ItalicFont=*-Italic, BoldItalicFont=*-SemiboldItalic]
\setsansfont{FiraSans}[Extension=.otf, UprightFont=*-Regular, BoldFont=*-SemiBold,
  ItalicFont=*-Italic, BoldItalicFont=*-SemiBoldItalic, Scale=MatchLowercase]
\setmonofont{FiraMono}[Extension=.otf, UprightFont=*-Regular, BoldFont=*-Medium, AutoFakeSlant=0.15, Scale=0.85]
\setmathfont{LibertinusMath-Regular.otf}

\usepackage{microtype}
\usepackage{enumitem}
\usepackage{multicol}
\setlist{itemsep=2pt, topsep=4pt}
\usepackage{caption}
\captionsetup{font=small, labelfont={sf,bf,color=pdsazul}, labelsep=period}
\usepackage{subcaption}
\usepackage[most]{tcolorbox}
\usepackage{environ}
\usepackage{scrlayer-scrpage}
\usepackage[hidelinks, unicode]{hyperref}
\usepackage[spanish, noabbrev]{cleveref}

% Títulos en Fira Sans con el azul de la asignatura
\setkomafont{disposition}{\sffamily\bfseries\color{pdsazul}}
\setkomafont{descriptionlabel}{\sffamily\bfseries}
\setlength{\parindent}{0pt}
\setlength{\parskip}{0.55\baselineskip plus 2pt}
\linespread{1.05}

% ---------------------------------------------------------------------
%  Datos del tema y cabeceras
% ---------------------------------------------------------------------
\newcommand{\numtema}{}
\newcommand{\nombretema}{}
\newcommand{\tema}[2]{\renewcommand{\numtema}{#1}\renewcommand{\nombretema}{#2}}

\pagestyle{scrheadings}
\clearpairofpagestyles
\ihead{\sffamily\small\color{pdsgris}\asignatura}
\ohead{\sffamily\small\color{pdsgris}Tema~\numtema. \nombretema}
\ifoot{\sffamily\small\color{pdsgris}\grado}
\ofoot{\sffamily\small\pagemark}
\setkomafont{pagehead}{\normalfont}
\setkomafont{pagefoot}{\normalfont}

% \cabecera{Apuntes|Ejercicios|...} : bloque de título de la primera página
\newcommand{\cabecera}[1]{%
  \thispagestyle{plain.scrheadings}%
  \noindent
  \begin{minipage}[c]{2.2cm}
    \includegraphics[width=1.9cm]{../../comun/ucm.png}
  \end{minipage}%
  \begin{minipage}[c]{\dimexpr\linewidth-2.2cm}
    \sffamily
    {\small\color{pdsgris}\grado\par \asignatura\ · \curso\par}
    \vspace{4pt}
    {\Large\bfseries\color{pdsazul}#1\par}
    {\large\bfseries Tema~\numtema. \nombretema\par}
    \vspace{2pt}
    {\small\color{pdsgris}\autores\par}
  \end{minipage}\par
  \vspace{6pt}{\color{pdsazul}\rule{\linewidth}{0.8pt}}\par\vspace{4pt}}

% ---------------------------------------------------------------------
%  Cajas: definición, propiedad, ejemplo, nota, «en código»
% ---------------------------------------------------------------------
\tcbset{pdsbase/.style={enhanced, breakable, boxrule=0pt, arc=1.5pt, left=8pt, right=8pt,
    top=5pt, bottom=5pt, fonttitle=\sffamily\bfseries\small, coltitle=white,
    before skip=8pt, after skip=8pt}}
\newtcbtheorem[number within=section]{definicion}{Definición}%
  {pdsbase, colback=pdsfondo, colframe=pdsazul, borderline west={2.5pt}{0pt}{pdsazul}}{def}
\newtcbtheorem[use counter from=definicion]{propiedad}{Propiedad}%
  {pdsbase, colback=pdsfondo, colframe=pdsazul!75, borderline west={2.5pt}{0pt}{pdsazul!75}}{prop}
\newtcbtheorem[use counter from=definicion]{ejemplo}{Ejemplo}%
  {pdsbase, colback=pdsverde!5, colframe=pdsverde, borderline west={2.5pt}{0pt}{pdsverde}}{ej}
\newtcolorbox{nota}[1][Nota]{pdsbase, colback=pdsnaranja!7, colframe=pdsnaranja,
  title={#1}, borderline west={2.5pt}{0pt}{pdsnaranja}}
\newtcolorbox{importante}[1][Importante]{pdsbase, colback=pdsrojo!5, colframe=pdsrojo,
  title={#1}, borderline west={2.5pt}{0pt}{pdsrojo}}
\newtcolorbox{encodigo}[1][En código]{pdsbase, colback=white, colframe=pdsgris,
  title={#1}, borderline west={2.5pt}{0pt}{pdsgris}}

\lstset{basicstyle=\ttfamily\small}
\renewcommand{\codigoetiqueta}[1]{{\sffamily\small\bfseries\color{pdsgris}#1}}

% ---------------------------------------------------------------------
%  Hojas de ejercicios: \begin{ejercicio} ... \end{ejercicio} y
%  \begin{solucion} ... \end{solucion} (solo se imprime si se define
%  \SOLUCIONES antes de cargar el documento; ver *-Soluciones.tex)
% ---------------------------------------------------------------------
\newif\ifsoluciones
\ifdefined\SOLUCIONES \solucionestrue \fi
\newcounter{ejercicio}
\newenvironment{ejercicio}[1][]{%
  \refstepcounter{ejercicio}\par\vspace{6pt}%
  \noindent{\sffamily\bfseries\color{pdsazul}Ejercicio~\theejercicio}%
  \if\relax\detokenize{#1}\relax\else\ {\sffamily\color{pdsgris}(#1)}\fi\quad}%
  {\par}
\NewEnviron{solucion}{%
  \ifsoluciones
    \begin{tcolorbox}[pdsbase, colback=pdsverde!4, colframe=pdsverde,
      borderline west={2.5pt}{0pt}{pdsverde}, title={Solución}]
      \BODY
    \end{tcolorbox}
  \fi}
% \begin{incisos}[columnas] ... \end{incisos}: apartados a), b)... de un
% enunciado repartidos en columnas (3 por defecto)
% (con una sola columna no se usa multicols, que necesita al menos dos)
\newenvironment{incisos}[1][3]{\par\vspace{-2pt}%
  \ifnum#1>1 \def\finincisos{\end{multicols}}\begin{multicols}{#1}\raggedcolumns\raggedright
  \else\def\finincisos{}\fi
  \begin{enumerate}[label=\alph*), itemsep=3pt, topsep=0pt]}%
  {\end{enumerate}\finincisos\vspace{-4pt}}
\newcommand{\herramienta}[1]{\emph{Herramienta interactiva:} \texttt{herramientas/#1}}

%     \tema{1}{Introducción: señales y sistemas}
%     \begin{document} \cabecera{Apuntes} ... \end{document}
%  Compilar con:  latexmk -xelatex  (o «make» desde teoría/)
% =====================================================================

% =====================================================================
%  pds-comun.tex — Paquetes, colores y estilos de código comunes a
%  transparencias, apuntes y hojas de ejercicios. Compilar con XeLaTeX.
%  No cargar directamente: lo cargan pds-beamer.tex y pds-apuntes.tex.
% =====================================================================

\usepackage{polyglossia}
\setmainlanguage{spanish}
\usepackage{mathtools}
\usepackage{graphicx}
\usepackage{booktabs}
\usepackage{multirow}
\usepackage{array}
\usepackage{xcolor}
\usepackage{tikz}
\usetikzlibrary{positioning,arrows.meta,calc,shapes.geometric,decorations.pathreplacing}
\usepackage{listings}
\usepackage{siunitx}
\sisetup{output-decimal-marker={,}, group-separator={\,}, per-mode=symbol}

% Notas de erratas de versiones anteriores: no se imprimen salvo que se
% defina \ERRATAS antes de cargar el estilo (make ERRATAS=1)
\newif\iferratas
\ifdefined\ERRATAS \erratastrue \fi

% Ruta común (logo y figuras compartidas)
\providecommand{\rutacomun}{../../comun/}
\input{\rutacomun datos.tex}

% ---------------------------------------------------------------------
%  Paleta (la misma que usan las figuras generadas con Python,
%  ver comun/pds_estilo.py)
% ---------------------------------------------------------------------
\definecolor{pdsazul}{HTML}{2B3A8C}
\definecolor{pdsrojo}{HTML}{C0392B}
\definecolor{pdsverde}{HTML}{1E7B4F}
\definecolor{pdsnaranja}{HTML}{D9822B}
\definecolor{pdsgris}{HTML}{5A5A66}
\definecolor{pdsfondo}{HTML}{F3F4FA}
\definecolor{pdscodigo}{HTML}{F7F7F2}

% ---------------------------------------------------------------------
%  Diagramas de bloques (TikZ)
% ---------------------------------------------------------------------
\tikzset{
  bloque/.style={rectangle, rounded corners=2pt, draw=pdsazul!80, fill=pdsfondo,
                 very thick, minimum height=8mm, minimum width=14mm, align=center},
  sumador/.style={circle, draw=pdsazul!80, fill=pdsfondo, thick, inner sep=1.5pt},
  flecha/.style={-{Stealth[length=2.2mm]}, thick, pdsgris},
}

% ---------------------------------------------------------------------
%  Código: Python y Matlab en paralelo
% ---------------------------------------------------------------------
\newcommand{\pdscodesize}{\footnotesize}
\lstdefinestyle{pds}{
  basicstyle=\ttfamily\pdscodesize,
  keywordstyle=\color{pdsazul}\bfseries,
  commentstyle=\color{pdsverde}\itshape,
  stringstyle=\color{pdsrojo},
  numbers=none,
  backgroundcolor=\color{pdscodigo},
  frame=single, rulecolor=\color{pdsgris!40}, framerule=0.4pt,
  xleftmargin=2pt, xrightmargin=2pt, aboveskip=2pt, belowskip=2pt,
  columns=fullflexible, keepspaces=true,
  showstringspaces=false, breaklines=true, tabsize=4,
  upquote=true,
  % Con XeLaTeX, listings desordena los caracteres no ASCII junto a
  % dígitos («2π» sale «π2»): se sustituyen por su versión matemática.
  literate={π}{{$\pi$}}1 {δ}{{$\delta$}}1 {ω}{{$\omega$}}1 {Ω}{{$\Omega$}}1
           {≈}{{$\approx$}}1 {µ}{{$\mu$}}1 {·}{{$\cdot$}}1 {α}{{$\alpha$}}1
           {á}{{á}}1 {é}{{é}}1 {í}{{í}}1 {ó}{{ó}}1 {ú}{{ú}}1 {ñ}{{ñ}}1 {¿}{{¿}}1,
}
\lstdefinestyle{python}{style=pds, language=Python,
  morekeywords={as,True,False,None}}
\lstdefinestyle{matlab}{style=pds, language=Matlab,
  morekeywords={function,end,true,false}}

% \codigo{<ruta sin extensión>} muestra el mismo ejemplo en Python
% y en Matlab, uno al lado del otro. Los ficheros viven en
% ../codigo/python/<nombre>.py y ../codigo/matlab/<nombre>.m y se
% ejecutan en las pruebas (make prueba-codigo).
\newcommand{\rutacodigo}{../codigo/}
\newcommand{\codigoetiqueta}[1]{{\sffamily\scriptsize\bfseries\color{pdsgris}#1}}
\newcommand{\codigo}[2][]{%
  \par\noindent
  \begin{minipage}[t]{0.495\linewidth}
    \codigoetiqueta{Python}\par\vspace{1pt}
    \lstinputlisting[style=python,#1]{\rutacodigo python/#2.py}
  \end{minipage}\hfill
  \begin{minipage}[t]{0.495\linewidth}
    \codigoetiqueta{Matlab}\par\vspace{1pt}
    \lstinputlisting[style=matlab,#1]{\rutacodigo matlab/#2.m}
  \end{minipage}\par}

\input{\rutacomun notacion.tex}


\usepackage[a4paper, top=2.6cm, bottom=2.6cm, left=2.5cm, right=2.5cm,
            headheight=1.1cm]{geometry}
\usepackage{fontspec}
\usepackage{unicode-math}
\setmainfont{LibertinusSerif}[Extension=.otf, UprightFont=*-Regular, BoldFont=*-Semibold,
  ItalicFont=*-Italic, BoldItalicFont=*-SemiboldItalic]
\setsansfont{FiraSans}[Extension=.otf, UprightFont=*-Regular, BoldFont=*-SemiBold,
  ItalicFont=*-Italic, BoldItalicFont=*-SemiBoldItalic, Scale=MatchLowercase]
\setmonofont{FiraMono}[Extension=.otf, UprightFont=*-Regular, BoldFont=*-Medium, AutoFakeSlant=0.15, Scale=0.85]
\setmathfont{LibertinusMath-Regular.otf}

\usepackage{microtype}
\usepackage{enumitem}
\usepackage{multicol}
\setlist{itemsep=2pt, topsep=4pt}
\usepackage{caption}
\captionsetup{font=small, labelfont={sf,bf,color=pdsazul}, labelsep=period}
\usepackage{subcaption}
\usepackage[most]{tcolorbox}
\usepackage{environ}
\usepackage{scrlayer-scrpage}
\usepackage[hidelinks, unicode]{hyperref}
\usepackage[spanish, noabbrev]{cleveref}

% Títulos en Fira Sans con el azul de la asignatura
\setkomafont{disposition}{\sffamily\bfseries\color{pdsazul}}
\setkomafont{descriptionlabel}{\sffamily\bfseries}
\setlength{\parindent}{0pt}
\setlength{\parskip}{0.55\baselineskip plus 2pt}
\linespread{1.05}

% ---------------------------------------------------------------------
%  Datos del tema y cabeceras
% ---------------------------------------------------------------------
\newcommand{\numtema}{}
\newcommand{\nombretema}{}
\newcommand{\tema}[2]{\renewcommand{\numtema}{#1}\renewcommand{\nombretema}{#2}}

\pagestyle{scrheadings}
\clearpairofpagestyles
\ihead{\sffamily\small\color{pdsgris}\asignatura}
\ohead{\sffamily\small\color{pdsgris}Tema~\numtema. \nombretema}
\ifoot{\sffamily\small\color{pdsgris}\grado}
\ofoot{\sffamily\small\pagemark}
\setkomafont{pagehead}{\normalfont}
\setkomafont{pagefoot}{\normalfont}

% \cabecera{Apuntes|Ejercicios|...} : bloque de título de la primera página
\newcommand{\cabecera}[1]{%
  \thispagestyle{plain.scrheadings}%
  \noindent
  \begin{minipage}[c]{2.2cm}
    \includegraphics[width=1.9cm]{../../comun/ucm.png}
  \end{minipage}%
  \begin{minipage}[c]{\dimexpr\linewidth-2.2cm}
    \sffamily
    {\small\color{pdsgris}\grado\par \asignatura\ · \curso\par}
    \vspace{4pt}
    {\Large\bfseries\color{pdsazul}#1\par}
    {\large\bfseries Tema~\numtema. \nombretema\par}
    \vspace{2pt}
    {\small\color{pdsgris}\autores\par}
  \end{minipage}\par
  \vspace{6pt}{\color{pdsazul}\rule{\linewidth}{0.8pt}}\par\vspace{4pt}}

% ---------------------------------------------------------------------
%  Cajas: definición, propiedad, ejemplo, nota, «en código»
% ---------------------------------------------------------------------
\tcbset{pdsbase/.style={enhanced, breakable, boxrule=0pt, arc=1.5pt, left=8pt, right=8pt,
    top=5pt, bottom=5pt, fonttitle=\sffamily\bfseries\small, coltitle=white,
    before skip=8pt, after skip=8pt}}
\newtcbtheorem[number within=section]{definicion}{Definición}%
  {pdsbase, colback=pdsfondo, colframe=pdsazul, borderline west={2.5pt}{0pt}{pdsazul}}{def}
\newtcbtheorem[use counter from=definicion]{propiedad}{Propiedad}%
  {pdsbase, colback=pdsfondo, colframe=pdsazul!75, borderline west={2.5pt}{0pt}{pdsazul!75}}{prop}
\newtcbtheorem[use counter from=definicion]{ejemplo}{Ejemplo}%
  {pdsbase, colback=pdsverde!5, colframe=pdsverde, borderline west={2.5pt}{0pt}{pdsverde}}{ej}
\newtcolorbox{nota}[1][Nota]{pdsbase, colback=pdsnaranja!7, colframe=pdsnaranja,
  title={#1}, borderline west={2.5pt}{0pt}{pdsnaranja}}
\newtcolorbox{importante}[1][Importante]{pdsbase, colback=pdsrojo!5, colframe=pdsrojo,
  title={#1}, borderline west={2.5pt}{0pt}{pdsrojo}}
\newtcolorbox{encodigo}[1][En código]{pdsbase, colback=white, colframe=pdsgris,
  title={#1}, borderline west={2.5pt}{0pt}{pdsgris}}

\lstset{basicstyle=\ttfamily\small}
\renewcommand{\codigoetiqueta}[1]{{\sffamily\small\bfseries\color{pdsgris}#1}}

% ---------------------------------------------------------------------
%  Hojas de ejercicios: \begin{ejercicio} ... \end{ejercicio} y
%  \begin{solucion} ... \end{solucion} (solo se imprime si se define
%  \SOLUCIONES antes de cargar el documento; ver *-Soluciones.tex)
% ---------------------------------------------------------------------
\newif\ifsoluciones
\ifdefined\SOLUCIONES \solucionestrue \fi
\newcounter{ejercicio}
\newenvironment{ejercicio}[1][]{%
  \refstepcounter{ejercicio}\par\vspace{6pt}%
  \noindent{\sffamily\bfseries\color{pdsazul}Ejercicio~\theejercicio}%
  \if\relax\detokenize{#1}\relax\else\ {\sffamily\color{pdsgris}(#1)}\fi\quad}%
  {\par}
\NewEnviron{solucion}{%
  \ifsoluciones
    \begin{tcolorbox}[pdsbase, colback=pdsverde!4, colframe=pdsverde,
      borderline west={2.5pt}{0pt}{pdsverde}, title={Solución}]
      \BODY
    \end{tcolorbox}
  \fi}
% \begin{incisos}[columnas] ... \end{incisos}: apartados a), b)... de un
% enunciado repartidos en columnas (3 por defecto)
% (con una sola columna no se usa multicols, que necesita al menos dos)
\newenvironment{incisos}[1][3]{\par\vspace{-2pt}%
  \ifnum#1>1 \def\finincisos{\end{multicols}}\begin{multicols}{#1}\raggedcolumns\raggedright
  \else\def\finincisos{}\fi
  \begin{enumerate}[label=\alph*), itemsep=3pt, topsep=0pt]}%
  {\end{enumerate}\finincisos\vspace{-4pt}}
\newcommand{\herramienta}[1]{\emph{Herramienta interactiva:} \texttt{herramientas/#1}}

%     \tema{1}{Introducción: señales y sistemas}
%     \begin{document} \cabecera{Apuntes} ... \end{document}
%  Compilar con:  latexmk -xelatex  (o «make» desde teoría/)
% =====================================================================

% =====================================================================
%  pds-comun.tex — Paquetes, colores y estilos de código comunes a
%  transparencias, apuntes y hojas de ejercicios. Compilar con XeLaTeX.
%  No cargar directamente: lo cargan pds-beamer.tex y pds-apuntes.tex.
% =====================================================================

\usepackage{polyglossia}
\setmainlanguage{spanish}
\usepackage{mathtools}
\usepackage{graphicx}
\usepackage{booktabs}
\usepackage{multirow}
\usepackage{array}
\usepackage{xcolor}
\usepackage{tikz}
\usetikzlibrary{positioning,arrows.meta,calc,shapes.geometric,decorations.pathreplacing}
\usepackage{listings}
\usepackage{siunitx}
\sisetup{output-decimal-marker={,}, group-separator={\,}, per-mode=symbol}

% Notas de erratas de versiones anteriores: no se imprimen salvo que se
% defina \ERRATAS antes de cargar el estilo (make ERRATAS=1)
\newif\iferratas
\ifdefined\ERRATAS \erratastrue \fi

% Ruta común (logo y figuras compartidas)
\providecommand{\rutacomun}{../../comun/}
\input{\rutacomun datos.tex}

% ---------------------------------------------------------------------
%  Paleta (la misma que usan las figuras generadas con Python,
%  ver comun/pds_estilo.py)
% ---------------------------------------------------------------------
\definecolor{pdsazul}{HTML}{2B3A8C}
\definecolor{pdsrojo}{HTML}{C0392B}
\definecolor{pdsverde}{HTML}{1E7B4F}
\definecolor{pdsnaranja}{HTML}{D9822B}
\definecolor{pdsgris}{HTML}{5A5A66}
\definecolor{pdsfondo}{HTML}{F3F4FA}
\definecolor{pdscodigo}{HTML}{F7F7F2}

% ---------------------------------------------------------------------
%  Diagramas de bloques (TikZ)
% ---------------------------------------------------------------------
\tikzset{
  bloque/.style={rectangle, rounded corners=2pt, draw=pdsazul!80, fill=pdsfondo,
                 very thick, minimum height=8mm, minimum width=14mm, align=center},
  sumador/.style={circle, draw=pdsazul!80, fill=pdsfondo, thick, inner sep=1.5pt},
  flecha/.style={-{Stealth[length=2.2mm]}, thick, pdsgris},
}

% ---------------------------------------------------------------------
%  Código: Python y Matlab en paralelo
% ---------------------------------------------------------------------
\newcommand{\pdscodesize}{\footnotesize}
\lstdefinestyle{pds}{
  basicstyle=\ttfamily\pdscodesize,
  keywordstyle=\color{pdsazul}\bfseries,
  commentstyle=\color{pdsverde}\itshape,
  stringstyle=\color{pdsrojo},
  numbers=none,
  backgroundcolor=\color{pdscodigo},
  frame=single, rulecolor=\color{pdsgris!40}, framerule=0.4pt,
  xleftmargin=2pt, xrightmargin=2pt, aboveskip=2pt, belowskip=2pt,
  columns=fullflexible, keepspaces=true,
  showstringspaces=false, breaklines=true, tabsize=4,
  upquote=true,
  % Con XeLaTeX, listings desordena los caracteres no ASCII junto a
  % dígitos («2π» sale «π2»): se sustituyen por su versión matemática.
  literate={π}{{$\pi$}}1 {δ}{{$\delta$}}1 {ω}{{$\omega$}}1 {Ω}{{$\Omega$}}1
           {≈}{{$\approx$}}1 {µ}{{$\mu$}}1 {·}{{$\cdot$}}1 {α}{{$\alpha$}}1
           {á}{{á}}1 {é}{{é}}1 {í}{{í}}1 {ó}{{ó}}1 {ú}{{ú}}1 {ñ}{{ñ}}1 {¿}{{¿}}1,
}
\lstdefinestyle{python}{style=pds, language=Python,
  morekeywords={as,True,False,None}}
\lstdefinestyle{matlab}{style=pds, language=Matlab,
  morekeywords={function,end,true,false}}

% \codigo{<ruta sin extensión>} muestra el mismo ejemplo en Python
% y en Matlab, uno al lado del otro. Los ficheros viven en
% ../codigo/python/<nombre>.py y ../codigo/matlab/<nombre>.m y se
% ejecutan en las pruebas (make prueba-codigo).
\newcommand{\rutacodigo}{../codigo/}
\newcommand{\codigoetiqueta}[1]{{\sffamily\scriptsize\bfseries\color{pdsgris}#1}}
\newcommand{\codigo}[2][]{%
  \par\noindent
  \begin{minipage}[t]{0.495\linewidth}
    \codigoetiqueta{Python}\par\vspace{1pt}
    \lstinputlisting[style=python,#1]{\rutacodigo python/#2.py}
  \end{minipage}\hfill
  \begin{minipage}[t]{0.495\linewidth}
    \codigoetiqueta{Matlab}\par\vspace{1pt}
    \lstinputlisting[style=matlab,#1]{\rutacodigo matlab/#2.m}
  \end{minipage}\par}

\input{\rutacomun notacion.tex}


\usepackage[a4paper, top=2.6cm, bottom=2.6cm, left=2.5cm, right=2.5cm,
            headheight=1.1cm]{geometry}
\usepackage{fontspec}
\usepackage{unicode-math}
\setmainfont{LibertinusSerif}[Extension=.otf, UprightFont=*-Regular, BoldFont=*-Semibold,
  ItalicFont=*-Italic, BoldItalicFont=*-SemiboldItalic]
\setsansfont{FiraSans}[Extension=.otf, UprightFont=*-Regular, BoldFont=*-SemiBold,
  ItalicFont=*-Italic, BoldItalicFont=*-SemiBoldItalic, Scale=MatchLowercase]
\setmonofont{FiraMono}[Extension=.otf, UprightFont=*-Regular, BoldFont=*-Medium, AutoFakeSlant=0.15, Scale=0.85]
\setmathfont{LibertinusMath-Regular.otf}

\usepackage{microtype}
\usepackage{enumitem}
\usepackage{multicol}
\setlist{itemsep=2pt, topsep=4pt}
\usepackage{caption}
\captionsetup{font=small, labelfont={sf,bf,color=pdsazul}, labelsep=period}
\usepackage{subcaption}
\usepackage[most]{tcolorbox}
\usepackage{environ}
\usepackage{scrlayer-scrpage}
\usepackage[hidelinks, unicode]{hyperref}
\usepackage[spanish, noabbrev]{cleveref}

% Títulos en Fira Sans con el azul de la asignatura
\setkomafont{disposition}{\sffamily\bfseries\color{pdsazul}}
\setkomafont{descriptionlabel}{\sffamily\bfseries}
\setlength{\parindent}{0pt}
\setlength{\parskip}{0.55\baselineskip plus 2pt}
\linespread{1.05}

% ---------------------------------------------------------------------
%  Datos del tema y cabeceras
% ---------------------------------------------------------------------
\newcommand{\numtema}{}
\newcommand{\nombretema}{}
\newcommand{\tema}[2]{\renewcommand{\numtema}{#1}\renewcommand{\nombretema}{#2}}

\pagestyle{scrheadings}
\clearpairofpagestyles
\ihead{\sffamily\small\color{pdsgris}\asignatura}
\ohead{\sffamily\small\color{pdsgris}Tema~\numtema. \nombretema}
\ifoot{\sffamily\small\color{pdsgris}\grado}
\ofoot{\sffamily\small\pagemark}
\setkomafont{pagehead}{\normalfont}
\setkomafont{pagefoot}{\normalfont}

% \cabecera{Apuntes|Ejercicios|...} : bloque de título de la primera página
\newcommand{\cabecera}[1]{%
  \thispagestyle{plain.scrheadings}%
  \noindent
  \begin{minipage}[c]{2.2cm}
    \includegraphics[width=1.9cm]{../../comun/ucm.png}
  \end{minipage}%
  \begin{minipage}[c]{\dimexpr\linewidth-2.2cm}
    \sffamily
    {\small\color{pdsgris}\grado\par \asignatura\ · \curso\par}
    \vspace{4pt}
    {\Large\bfseries\color{pdsazul}#1\par}
    {\large\bfseries Tema~\numtema. \nombretema\par}
    \vspace{2pt}
    {\small\color{pdsgris}\autores\par}
  \end{minipage}\par
  \vspace{6pt}{\color{pdsazul}\rule{\linewidth}{0.8pt}}\par\vspace{4pt}}

% ---------------------------------------------------------------------
%  Cajas: definición, propiedad, ejemplo, nota, «en código»
% ---------------------------------------------------------------------
\tcbset{pdsbase/.style={enhanced, breakable, boxrule=0pt, arc=1.5pt, left=8pt, right=8pt,
    top=5pt, bottom=5pt, fonttitle=\sffamily\bfseries\small, coltitle=white,
    before skip=8pt, after skip=8pt}}
\newtcbtheorem[number within=section]{definicion}{Definición}%
  {pdsbase, colback=pdsfondo, colframe=pdsazul, borderline west={2.5pt}{0pt}{pdsazul}}{def}
\newtcbtheorem[use counter from=definicion]{propiedad}{Propiedad}%
  {pdsbase, colback=pdsfondo, colframe=pdsazul!75, borderline west={2.5pt}{0pt}{pdsazul!75}}{prop}
\newtcbtheorem[use counter from=definicion]{ejemplo}{Ejemplo}%
  {pdsbase, colback=pdsverde!5, colframe=pdsverde, borderline west={2.5pt}{0pt}{pdsverde}}{ej}
\newtcolorbox{nota}[1][Nota]{pdsbase, colback=pdsnaranja!7, colframe=pdsnaranja,
  title={#1}, borderline west={2.5pt}{0pt}{pdsnaranja}}
\newtcolorbox{importante}[1][Importante]{pdsbase, colback=pdsrojo!5, colframe=pdsrojo,
  title={#1}, borderline west={2.5pt}{0pt}{pdsrojo}}
\newtcolorbox{encodigo}[1][En código]{pdsbase, colback=white, colframe=pdsgris,
  title={#1}, borderline west={2.5pt}{0pt}{pdsgris}}

\lstset{basicstyle=\ttfamily\small}
\renewcommand{\codigoetiqueta}[1]{{\sffamily\small\bfseries\color{pdsgris}#1}}

% ---------------------------------------------------------------------
%  Hojas de ejercicios: \begin{ejercicio} ... \end{ejercicio} y
%  \begin{solucion} ... \end{solucion} (solo se imprime si se define
%  \SOLUCIONES antes de cargar el documento; ver *-Soluciones.tex)
% ---------------------------------------------------------------------
\newif\ifsoluciones
\ifdefined\SOLUCIONES \solucionestrue \fi
\newcounter{ejercicio}
\newenvironment{ejercicio}[1][]{%
  \refstepcounter{ejercicio}\par\vspace{6pt}%
  \noindent{\sffamily\bfseries\color{pdsazul}Ejercicio~\theejercicio}%
  \if\relax\detokenize{#1}\relax\else\ {\sffamily\color{pdsgris}(#1)}\fi\quad}%
  {\par}
\NewEnviron{solucion}{%
  \ifsoluciones
    \begin{tcolorbox}[pdsbase, colback=pdsverde!4, colframe=pdsverde,
      borderline west={2.5pt}{0pt}{pdsverde}, title={Solución}]
      \BODY
    \end{tcolorbox}
  \fi}
% \begin{incisos}[columnas] ... \end{incisos}: apartados a), b)... de un
% enunciado repartidos en columnas (3 por defecto)
% (con una sola columna no se usa multicols, que necesita al menos dos)
\newenvironment{incisos}[1][3]{\par\vspace{-2pt}%
  \ifnum#1>1 \def\finincisos{\end{multicols}}\begin{multicols}{#1}\raggedcolumns\raggedright
  \else\def\finincisos{}\fi
  \begin{enumerate}[label=\alph*), itemsep=3pt, topsep=0pt]}%
  {\end{enumerate}\finincisos\vspace{-4pt}}
\newcommand{\herramienta}[1]{\emph{Herramienta interactiva:} \texttt{herramientas/#1}}

\tema{3}{Transformada discreta de Fourier}

\begin{document}
\ifsoluciones\cabecera{Ejercicios resueltos}\else\cabecera{Ejercicios}\fi

Los ejercicios 1–4 corresponden a la parte~1 del tema (de la DTFT a la DFT), los ejercicios 5–9 a la parte~2 (propiedades y convolución circular), los ejercicios 10–13 a la parte~3 (análisis espectral) y los ejercicios 14–16 a la parte~4 (FFT y convolución rápida). Se usa $W_N=\e^{-j2\pi/N}$ y $((n))_N=n\bmod N$. Las secuencias se dan como listas que empiezan en $n=0$. Muchos resultados se pueden comprobar con \texttt{fft} y con los ejemplos de \texttt{teoría/T3/codigo/}.

% =====================================================================
\section*{De la DTFT a la DFT}

\begin{ejercicio}[DFT a mano]
  Calcula la DFT de 4 puntos de $x=(1,2,0,1)$ con la matriz $\mathbf W_4$, y comprueba el resultado aplicando la IDFT.
\end{ejercicio}
\begin{solucion}
  Con $W_4=-j$:
  \[
    \begin{pmatrix}X[0]\\X[1]\\X[2]\\X[3]\end{pmatrix}=
    \begin{pmatrix}1&1&1&1\\1&-j&-1&j\\1&-1&1&-1\\1&j&-1&-j\end{pmatrix}
    \begin{pmatrix}1\\2\\0\\1\end{pmatrix}=
    \begin{pmatrix}4\\1-j\\-2\\1+j\end{pmatrix}.
  \]
  $X[3]=X^*[1]$ porque $x$ es real. La IDFT, $\frac14\mathbf W_4^*\mathbf X$, da por ejemplo $x[0]=\frac14(4+1-j-2+1+j)=1$ y $x[1]=\frac14\big(4+j(1-j)+2-j(1+j)\big)=\frac14(4+1+j+2+1-j)=2$.
\end{solucion}

\begin{ejercicio}[DFT de secuencias básicas]
  Calcula, sin hacer las sumas término a término, la DFT de $N=8$ puntos de:
  \begin{incisos}[3]
    \item $\delta[n]$
    \item $\delta[n-2]$
    \item $1$ para todo $n$
    \item $\cos(\pi n/2)$
    \item $\e^{j\pi n/4}$
    \item $(-1)^n$
  \end{incisos}
\end{ejercicio}
\begin{solucion}
  \begin{enumerate}[label=\alph*)]
    \item $X[k]=1$ para todo $k$.
    \item $X[k]=W_8^{2k}=(-j)^k=(1,-j,-1,j,1,-j,-1,j)$.
    \item $X[k]=8\,\delta[k]$: solo la componente continua.
    \item $\cos(\pi n/2)=\frac12\big(\e^{j\frac{2\pi}{8}2n}+\e^{-j\frac{2\pi}{8}2n}\big)$ da dos vueltas en 8 muestras: $X[2]=X[6]=4$ y el resto, cero.
    \item Da una vuelta: $X[k]=8\,\delta[k-1]$.
    \item $(-1)^n=\e^{j\frac{2\pi}{8}4n}$: $X[k]=8\,\delta[k-4]$, la frecuencia $\Omega=\pi$.
  \end{enumerate}
\end{solucion}

\begin{ejercicio}[aliasing temporal]
  La secuencia $x=(1,2,3,4,5,6)$ tiene DTFT $\Xd$. Se toman $N=4$ muestras equiespaciadas de $\Xd$ en $[0,2\pi)$ y se calcula su IDFT.
  \begin{incisos}[1]
    \item ¿Qué secuencia se obtiene?
    \item ¿Cuál es el menor $N$ que permite recuperar $x[n]$?
  \end{incisos}
\end{ejercicio}
\begin{solucion}
  \begin{enumerate}[label=\alph*)]
    \item Un periodo de la extensión periódica de periodo 4: $\tilde x[n]=x[n]+x[n+4]$, es decir, $(1+5,\ 2+6,\ 3,\ 4)=(6,8,3,4)$.
    \item $N\ge L=6$.
  \end{enumerate}
\end{solucion}

\begin{ejercicio}[DFT de una secuencia infinita]
  Sea $x[n]=0{,}5^nu[n]$.
  \begin{incisos}[1]
    \item Obtén $X[k]$, las $N=4$ muestras de su DTFT en $\Omega_k=2\pi k/4$.
    \item Calcula la IDFT de esas muestras y compárala con $x[n]$ para $n=0,\dots,3$.
    \item ¿Qué $N$ hace falta para que el error sea menor que $10^{-3}$?
  \end{incisos}
\end{ejercicio}
\begin{solucion}
  \begin{enumerate}[label=\alph*)]
    \item $X[k]=\dfrac{1}{1-0{,}5W_4^k}$: $X=(2;\ 0{,}8-0{,}4j;\ 0{,}667;\ 0{,}8+0{,}4j)$.
    \item $\tilde x[n]=\sum_mx[n+4m]=\frac{0{,}5^n}{1-0{,}5^4}=1{,}0667\cdot0{,}5^n$: $(1{,}067;\ 0{,}533;\ 0{,}267;\ 0{,}133)$, cada valor un $6{,}7$\,\% mayor que $x[n]$ por el aliasing temporal.
    \item El error relativo es $\frac{0{,}5^N}{1-0{,}5^N}\approx0{,}5^N$, y el absoluto, como mucho en $n=0$, lo mismo. $0{,}5^N<10^{-3}$ exige $N\ge10$.
  \end{enumerate}
\end{solucion}

% =====================================================================
\section*{Propiedades y convolución circular}

\begin{ejercicio}[simetría y Parseval]
  Una secuencia real de longitud 8 tiene $X[0]=10$, $X[1]=2+j$, $X[2]=-1-2j$, $X[3]=3$ y $X[4]=-2$.
  \begin{incisos}[1]
    \item Completa $X[5]$, $X[6]$ y $X[7]$.
    \item Calcula $x[0]$ y $\sum_nx[n]$ sin calcular la IDFT completa.
    \item Calcula la energía $\sum_n\abs{x[n]}^2$.
  \end{incisos}
\end{ejercicio}
\begin{solucion}
  \begin{enumerate}[label=\alph*)]
    \item $X[k]=X^*[8-k]$: $X[5]=X^*[3]=3$, $X[6]=X^*[2]=-1+2j$, $X[7]=X^*[1]=2-j$.
    \item $\sum_nx[n]=X[0]=10$. $x[0]=\frac18\sum_kX[k]=\frac18(10+2+j-1-2j+3-2+3-1+2j+2-j)=\frac{16}{8}=2$.
    \item Por Parseval, $\frac18\sum_k\abs{X[k]}^2=\frac18(100+5+5+9+4+9+5+5)=\frac{142}8=17{,}75$.
  \end{enumerate}
\end{solucion}

\begin{ejercicio}[desplazamiento y modulación]
  Para $x=(1,2,0,1)$, con la DFT del ejercicio~1, obtén sin calcular ninguna DFT nueva:
  \begin{incisos}[1]
    \item la DFT de $y[n]=x[((n-2))_4]$;
    \item la DFT de $z[n]=(-1)^nx[n]$;
    \item la secuencia cuya DFT es $X[((-k))_4]$.
  \end{incisos}
\end{ejercicio}
\begin{solucion}
  \begin{enumerate}[label=\alph*)]
    \item $y=(0,1,1,2)$ y $Y[k]=W_4^{2k}X[k]=(-1)^kX[k]=(4,\ -1+j,\ -2,\ -1-j)$.
    \item $(-1)^n=W_4^{-2n}$, así que $Z[k]=X[((k-2))_4]=(X[2],X[3],X[0],X[1])=(-2,\ 1+j,\ 4,\ 1-j)$; $z=(1,-2,0,-1)$.
    \item La inversión circular $x[((-n))_4]=(x[0],x[3],x[2],x[1])=(1,1,0,2)$. Como $x$ es real, $X[((-k))_4]=X^*[k]$.
  \end{enumerate}
\end{solucion}

\begin{ejercicio}[convolución circular a mano]
  Sean $x=(1,2,3,4)$ y $h=(1,1)$.
  \begin{incisos}[1]
    \item Calcula la convolución lineal $x\conv h$.
    \item Calcula la convolución circular de 4 puntos y compárala con la lineal.
    \item ¿Qué $N$ hace falta para que coincidan?
  \end{incisos}
\end{ejercicio}
\begin{solucion}
  \begin{enumerate}[label=\alph*)]
    \item $y[n]=x[n]+x[n-1]=(1,3,5,7,4)$.
    \item $(x\cconv4h)[n]=x[n]+x[((n-1))_4]$: $(1+4,\ 2+1,\ 3+2,\ 4+3)=(5,3,5,7)$. Es la lineal plegada con periodo 4: la muestra $y[4]=4$ se suma a $y[0]=1$.
    \item $N\ge L+M-1=5$.
  \end{enumerate}
\end{solucion}

\begin{ejercicio}[longitud de la DFT para filtrar]
  Una señal de $L=100$ muestras se filtra con un FIR de $M=30$ coeficientes usando la DFT.
  \begin{incisos}[1]
    \item ¿Cuál es el menor $N$ que da la convolución lineal? ¿Y la menor potencia de 2?
    \item Si por error se usa $N=128$, ¿qué muestras de la salida son incorrectas?
    \item Si se usa $N=100$, ¿qué muestras de la salida son correctas?
  \end{incisos}
\end{ejercicio}
\begin{solucion}
  \begin{enumerate}[label=\alph*)]
    \item La salida tiene $L+M-1=129$ muestras: $N\ge129$, y la menor potencia de 2 es 256.
    \item La circular de 128 puntos es $y[n]+y[n+128]$, y solo $y[128]$ cae fuera: se suma a $y[0]$. Solo $y[0]$ es incorrecta.
    \item Ahora $y_c[n]=y[n]+y[n+100]$ con $y[n+100]\ne0$ para $n+100\le128$, es decir, $n\le28$. Son correctas las muestras $n=29,\dots,99$: todas menos las $M-1=29$ primeras. Es el fundamento del solapamiento y almacenamiento.
  \end{enumerate}
\end{solucion}

\begin{ejercicio}[secuencias pares]
  \begin{incisos}[1]
    \item Calcula la DFT de 4 puntos de $x=(4,1,0,1)$ y explica por qué es real.
    \item Demuestra que si $x[n]$ es real y par circular ($x[n]=x[((-n))_N]$), su DFT es real y par.
  \end{incisos}
\end{ejercicio}
\begin{solucion}
  \begin{enumerate}[label=\alph*)]
    \item $X[k]=4+W_4^k+W_4^{3k}=4+2\cos(\pi k/2)$: $X=(6,4,2,4)$. Es real porque $x[1]=x[3]$: la secuencia es par circular.
    \item Por la inversión, $x[((-n))_N]\leftrightarrow X[((-k))_N]$, así que $X[k]=X[((-k))_N]$ (par). Por ser real, $X[k]=X^*[((-k))_N]$. Juntando las dos, $X[k]=X^*[k]$: es real.
  \end{enumerate}
\end{solucion}

% =====================================================================
\section*{Análisis espectral}

\begin{ejercicio}[eje de frecuencias]
  Se calcula la DFT de $N=512$ muestras de una señal muestreada a $f_s=8$\,kHz.
  \begin{incisos}[1]
    \item ¿Cuál es la separación entre muestras de la DFT y la duración del registro?
    \item ¿En qué muestras aparece un tono de 1\,kHz?
    \item ¿A qué frecuencia corresponde la muestra $k=300$?
  \end{incisos}
\end{ejercicio}
\begin{solucion}
  \begin{enumerate}[label=\alph*)]
    \item $\Delta f=8000/512=15{,}625$\,Hz; $T_{\text{obs}}=512/8000=64$\,ms $=1/\Delta f$.
    \item $k_0=1000/15{,}625=64$ y su imagen $N-k_0=448$.
    \item $k=300>N/2=256$: es una frecuencia negativa, $(300-512)\cdot15{,}625=-3312{,}5$\,Hz. Si la señal es real, $X[300]=X^*[212]$, que corresponde a $+3312{,}5$\,Hz.
  \end{enumerate}
\end{solucion}

\begin{ejercicio}[amplitudes con y sin ventana]
  La señal $x(t)=3\cos(2\pi\,1000\,t)+\cos(2\pi\,2500\,t+\pi/4)$ se muestrea a $f_s=8$\,kHz y se toman $N=64$ muestras.
  \begin{incisos}[1]
    \item ¿En qué muestras $k\le N/2$ aparecen los tonos? ¿Cuánto valen $\abs{X[k]}$ y $\arg X[k]$ en ellas?
    \item Repite el cálculo con la ventana de Hann periódica, $w[n]=\frac12-\frac12\cos(2\pi n/N)$. ¿Qué valen las muestras vecinas de cada pico?
  \end{incisos}
\end{ejercicio}
\begin{solucion}
  \begin{enumerate}[label=\alph*)]
    \item $\Delta f=125$\,Hz: los tonos caen exactamente en $k=8$ y $k=20$. $\abs{X[8]}=\frac{3\cdot64}2=96$ con fase 0, y $\abs{X[20]}=\frac{64}{2}=32$ con fase $\pi/4$. El resto de las muestras son nulas.
    \item $w[n]=\frac12-\frac14\e^{j2\pi n/N}-\frac14\e^{-j2\pi n/N}$, así que cada línea se convierte en tres: $\frac12$ en $k_0$ y $-\frac14$ en $k_0\pm1$, multiplicando la DFT sin ventana. Con $\sum w=N/2=32$: $X_w[8]=48$, $X_w[7]=X_w[9]=-24$; $\abs{X_w[20]}=16$ y $\abs{X_w[19]}=\abs{X_w[21]}=8$. La amplitud se recupera con $2\abs{X_w[k_0]}/\sum w$: $2\cdot48/32=3$.
  \end{enumerate}
\end{solucion}

\begin{ejercicio}[resolución]
  Se quieren distinguir dos tonos de 1000 y 1030\,Hz muestreados a $f_s=8$\,kHz.
  \begin{incisos}[1]
    \item ¿Cuántas muestras hacen falta, aproximadamente, con la ventana rectangular? ¿Y con la de Hann?
    \item ¿Se pueden distinguir con $N=128$ muestras rellenando con ceros hasta 4096? Justifícalo.
  \end{incisos}
\end{ejercicio}
\begin{solucion}
  \begin{enumerate}[label=\alph*)]
    \item Con la rectangular, $f_s/N<30$\,Hz: $N>267$, unos 33\,ms de señal. Con Hann, el lóbulo principal es el doble de ancho: $N>533$, unos 67\,ms.
    \item No. El relleno con ceros da más muestras de la DTFT de las mismas 128 muestras, cuyo lóbulo principal (semianchura $f_s/N=62{,}5$\,Hz) es mucho más ancho que la separación de los tonos. Se ve un único pico.
  \end{enumerate}
\end{solucion}

\begin{ejercicio}[potencia y densidad espectral]
  La señal $x[n]=2\cos(2\pi\,100\,n/f_s)+v[n]$, con $f_s=1$\,kHz y $v[n]$ ruido blanco de varianza $\sigma^2=0{,}01$, se analiza con $N=1000$ muestras.
  \begin{incisos}[1]
    \item ¿Cuánto vale la potencia por muestra, $\abs{X[k]/N}^2$, en el tono y, en promedio, en el ruido? Exprésalo en dB.
    \item ¿Y la densidad espectral estimada por el periodograma, en dB/Hz?
    \item ¿Qué cambia si se duplica $N$?
  \end{incisos}
\end{ejercicio}
\begin{solucion}
  \begin{enumerate}[label=\alph*)]
    \item El tono cae en $k_0=100$: $\abs{X[k_0]}=AN/2$ y la potencia por muestra es $A^2/4=1$, 0\,dB. Para el ruido, $E\abs{V[k]}^2=N\sigma^2$, así que la potencia media por muestra es $\sigma^2/N=10^{-5}$, $-50$\,dB.
    \item Ruido: $\sigma^2/f_s=10^{-5}$\,/Hz, $-50$\,dB/Hz. Tono: $\frac{(AN/2)^2}{Nf_s}=\frac{A^2N}{4f_s}=1$, 0\,dB/Hz.
    \item La potencia por muestra del tono sigue siendo 0\,dB y la del ruido baja a $-53$\,dB. En densidad, el ruido sigue en $-50$\,dB/Hz y el pico del tono sube a 3\,dB/Hz. En los dos casos el tono sobresale 3\,dB más: alargar el registro ayuda a detectar tonos en ruido. (Simulado promediando 200 realizaciones.)
  \end{enumerate}
\end{solucion}

% =====================================================================
\section*{FFT y convolución rápida}

\begin{ejercicio}[coste de la FFT]
  Un procesador hace un producto complejo cada nanosegundo. ¿Cuánto tarda en calcular una DFT de $N=4096$ puntos directamente y con una FFT en base 2? Ignora las sumas.
\end{ejercicio}
\begin{solucion}
  Directa: $N^2=1{,}68\cdot10^7$ productos, $16{,}8$\,ms. FFT: $\frac N2\log_2N=2048\cdot12=24\,576$ productos, $24{,}6$\,µs. Es unas 680 veces más rápida.
\end{solucion}

\begin{ejercicio}[FFT de cuatro puntos]
  Calcula la DFT de $x=(1,3,4,0)$ con el algoritmo de diezmado en el tiempo: separa las muestras pares e impares, calcula sus DFT de 2 puntos y combínalas con las mariposas. Dibuja el diagrama de flujo y cuenta los productos.
\end{ejercicio}
\begin{solucion}
  Pares $(1,4)$: $F=(1+4,\ 1-4)=(5,-3)$. Impares $(3,0)$: $G=(3,3)$. Con $W_4^0=1$ y $W_4^1=-j$:
  \begin{align*}
    X[0]&=F[0]+G[0]=8, & X[1]&=F[1]-jG[1]=-3-3j,\\
    X[2]&=F[0]-G[0]=2, & X[3]&=F[1]+jG[1]=-3+3j.
  \end{align*}
  Coincide con la DFT directa. Hay $\frac N2\log_2N=4$ productos, todos por $\pm1$ o $\pm j$, es decir, triviales; la DFT directa necesita 16.
\end{solucion}

\begin{ejercicio}[solapamiento y suma]
  Un sistema de audio a 48\,kHz filtra en tiempo real con un FIR de $M=101$ coeficientes por solapamiento y suma con FFT de $N=1024$ puntos.
  \begin{incisos}[1]
    \item ¿Qué longitud $L$ deben tener los bloques? ¿Cuántos bloques se procesan por segundo?
    \item ¿Cuántos productos complejos por muestra de salida necesita, si $H[k]$ se calcula una vez? Compáralo con la convolución directa.
    \item ¿Qué retardo introduce el procesado por bloques?
  \end{incisos}
\end{ejercicio}
\begin{solucion}
  \begin{enumerate}[label=\alph*)]
    \item $L=N-M+1=924$ muestras; $48\,000/924\approx52$ bloques por segundo.
    \item Por bloque: una FFT, $N$ productos y una IFFT, $2\cdot512\cdot10+1024=11\,264$ productos complejos, unos $12{,}2$ por muestra de salida (unos 49 reales). La convolución directa necesita 101 productos reales por muestra: la FFT es unas dos veces más barata, y la ventaja crece con $M$.
    \item Al menos un bloque, $924/48\,000\approx19$\,ms, más el tiempo de cálculo.
  \end{enumerate}
\end{solucion}

\end{document}
